\section{Market Generator Framework}\label{sec:generator-framework}
It is important to have standardized generation framework as
part of our efforts to create a unified benchmark \ref{contribute:3};
this section outines the details of the dataset used, preprocessing, and
 generation framework used by StonkBench for models to train on 
 historical market data.

\subsection{Data preprocessing}\label{subsec:data-preprocess}
    \subsubsection{Dataset Split}
    Given our stock price dataset, we perform a chronological
    ten-two-three year partition into training, validation, and testing
    datasets from 2010 to 2025, where $D = D^{\mathrm{train}}, D^{\mathrm{val}}, D^{\mathrm{test}}$ denotes the concatenation.
    This deviates from the training scheme proposed by TSGBench \cite{Ang23}; we wish to 
    hold out our testing time-horizon from the training set completely
    to evaluate model generalization capabilities.
    
    \subsubsection{Log Return Transformation}
    We then transform the price time series for both sets into a log return time series, 
    which is a common practice in financial modeling to stabilize variance 
    and make the data more suitable for analysis. The return \( \mathbf{R}_t \) at time \( t \) is calculated as
    \footnote{The log return operation is applied channel-wise to the time series.}
    \footnote{From now, on we refer to $D$ as the log return time series dataset and notate it as $\mathbf{R}_t$
    to avoid notation clutering, unless explicitly stated otherwise.}:
    \begin{equation}
    \mathbf{R}_t := \log\left(\frac{\mathbf{S}_t}{\mathbf{S}_{t-1}}\right) = \log(\mathbf{S}_t) - \log(\mathbf{S}_{t-1}),
    \end{equation}

    \subsubsection{Sliding Window Transformation}
    The deep learning models considered in §\ref{sub:deep-learning-methods} require a fixed-length data
    for training and generation. Therefore, to standardize input and data generation, while preserving temporal dependencies,
    we apply a sliding window transformation to the log return time series. Given a
    window length $L\in\mathbb{N}$, we applying a sliding window operation $W_L$ over the time
    series as follows.
    \begin{equation}
        W_L: (\mathbf{R}_t)_{t=1}^T \mapsto \{ \mathbf{w}_t \coloneq(\mathbf{R}_t,\mathbf{R}_{t+1},\dots, \mathbf{R}_{t+L}): 1\leq t \leq T-L\}
    \end{equation}
    We denote the resulting set of windows, for a given data partition from $D$, as $\mathcal{D}^i\coloneq W_L(D^i)$ for $i\in\{\mathrm{train}, \mathrm{val}, \mathrm{test}\}$.
    For the purpose of research, we expirment our generation training with windown lenghts $L\in\{ 21, 63, 126 ,252\}$ , represeneting
    general the number of trading days in a month, a quarter, half a year and a year, respectively
    
    \footnote{TSGBench determined the window length for deep learning training 
    and generation depending on the dataset by finding the maximum auto-correlation of that time
    series dataset \cite{Ang23}}.


\subsection{Market Generator Training and Generation}\label{subsec:generator-training}
    We train each market generator on the log-return training data to
    produce synthetic windows, each with shape $L\times I$. Formally,
    for each generator $G_g$, $g\in\mathcal{G}$, we generate $N$ synthetic
    windows and obtain the generated set
    $\hat{\mathcal{D}}_g = \{\hat{\mathbf{w}}_n: \hat{\mathbf{w}}_n \sim G_g(\mathbf{w} \mid D^{\mathrm{train}})\}_{n=1}^N$.
    Here, $\mathcal{G}$ is the index set of market generators considered in
    §\ref{sec:model-overview}, $\hat{\mathbf{w}}_n$ denotes a synthetic
    window of size $L\times I$, and $N$ is the number of synthetic windows
    generated for each model, which we set to $N=1000$.

    \subsubsection{Non-deep Learning Model Generation}
    Parametric and bootstrap generators are fit on $D^{\mathrm{train}}$ and then
    sample windows directly from the learned process or resampled data.
    \begin{align}
        \hat{\mathbf{S}}_{n_l} \sim G_g(\mathbf{S} \mid D^{\mathrm{train}}), \quad
        \hat{\mathbf{w}}_n \coloneq (\hat{\mathbf{S}}_{n_1},\dots, \hat{\mathbf{S}}_{n_L}), \quad
        \hat{\mathcal{D}}_g = \{\hat{\mathbf{w}}_n\}_{n=1}^N.
    \end{align}

    \subsubsection{Deep Learning Model Generation}
    Deep-learning frameworks in this paper (GAN, VAE, DDPM, and
    signature-based models) are trained on the sliding-window transform
    $W_L(D^{\mathrm{train}})$ (see §\ref{subsec:data-preprocess}). They then
    generate $N$ independent windows of size $L$, each with shape
    $L\times I$, directly from the learned window distribution:
    \begin{align}
        \hat{\mathbf{w}}_n \sim G_g(\mathbf{w} \mid W_LD^{\mathrm{train}}), \qquad
        \hat{\mathcal{D}}_g = \{\hat{\mathbf{w}}_n\}_{n=1}^N.
    \end{align}










